% Pacotes e declaração do documento
\documentclass{article}
\usepackage[utf8]{inputenc}
\usepackage{graphicx}
\usepackage{amsmath}
\usepackage{indentfirst}
\usepackage{amsthm}
\usepackage{fourier}


% Customizações
\title{\textbf{C-Bounded Powersort's time complexity and correctness analysis}}
\author{Théo Araújo Magalhães}
\date{April 5, 2025}
\DeclareMathAlphabet{\mathcal}{OMS}{cmsy}{m}{n}
\SetMathAlphabet{\mathcal}{bold}{OMS}{cmsy}{b}{n}
\newcommand{\bigO}{\mathcal{O}}
\DeclareUnicodeCharacter{2212}{-}

%Documento
\begin{document}
\maketitle
\section*{Definitions:}
\textbf{Runs:} Contiguous regions of the input that are already in sorted order. Let R$_0$, \dots , R$_{r-1}$ be the runs in the input, L$_i$ denotes the length of R$_i$.*

\textbf{Run-adaptive mergesort:} Any stable run-adaptive mergesort can be described as a merge tree T: an (ordered rooted) tree with leaves R$_0$, \dots, R$_{r-1}$ (appearing in that order from left to right). Each internal node v of T corresponds to the result of merging its children.*

\textbf{Merge costs:} M(v) is the total length of the resulting merge corresponding to the internal node v. The total merge cost of T is then:
\begin{equation}
M(T) = \sum_{v \in T}M(v).     
\end{equation}

\textbf{Run boundaries:} Denoted by B$_i$, it is the boundary between the (i − 1)st and ith run (i=1, \dots, r - 1). In binary merge trees, there is a one-to-one correspondence between the B$_i$ and the (inner/non-leaf) merge-tree nodes.*

\textbf{Node depth:} Let d$_i$ be the depth of leaf R$_i$ in T, where depth is the number of edges on the path to the root. Every element in run R$_i$ is counted exactly d$_i$ times in $\sum$$_v$M(v), so: 
\begin{equation}
        M = \sum^{r-1}_{i=0}d_i \cdot L_i*    
\end{equation}

\textbf{Node power:} Let $\alpha$$_0$, \dots, $\alpha$$_m$, $\sum$$\alpha$$_j$ = 1 be leaf probabilities. For 1 $\le$ j $\le$ m, let j be the internal node separating the (j − 1)st and jth leaf. The power of (the split at) node j is:
\begin{equation}
    P_j = min\{ \ell \in N: \lfloor{\frac{a}{2^{-\ell}}}\rfloor   < \lfloor{\frac{b}{2^{-\ell}}}\rfloor \}, \text{where } a = \sum_{i=0}^{j-1} \alpha_i - \frac{1}{2}\alpha_{j-1}, b = \sum_{i=0}^{j-1} \alpha_i + \frac{1}{2}\alpha_j.
\end{equation}

(Pj is the index of the first bit where the (binary) fractional parts of a and b differ.)

\textbf{Powersort:} The intuition behind Powersort is to imagine a complete binary tree T$_v$ sitting on top of the array, where each element A[i] of the input corresponds to a leaf of T$_v$ .T$_v$ corresponds to the merge tree of a non-adaptive two-way Mergesort (starting with individual elements). Now each run boundary B$_j$ also corresponds to an (inner) node w of this (virtual) k-ary tree TV , namely the lowest common ancestor of the two leafs adjacent to B$_j$.* [ We need to re-explain this part better so that we can talk about the concept of a “Powersort merge tree” or something of the sorts. Helps with the proof of our main proposition]  

\textbf{Forgetful K-Stack:} A forgetful stack is a stack that forgets its k-th element if it already has k elements in it and a new one is added. 

\textbf{C-Bounded Powersort:} The Powersort algorithm using a Forgetful C-Stack instead of the traditional stack with size logn + 1. 

\textbf{Processing a subtree of height h:} Let T be a merge tree, T’ one of its subtrees of size h and v the root of T’. Processing T’ is executing all the merges necessary to fill node v, given that it has length M(v). 

\section*{Proofs:} 

\textbf{Proposition:} Let T be a Powersort merge tree of height H. One pass of the C-Bounded Powersort algorithm is enough to process all subtrees of T with height of C-1.

\textit{Proof by induction:}  

\textbf{Base case - T of height 2:}

Let v be the root of T and v$_1$ and v$_2$ its two children. Therefore, T has two sub-trees T’ a T’’ of size one. For the left subtree T', v$_1$ is the root and it has only two leaf children: the runs R$_1$ and R$_2$. Given that h = 1, a Forgetful 2-Stack will work with 2-Bounded Powersort as follows: 

R$_1$  will be found and added to the stack, then, given that there is nothing to compare R$_1$  with, the algorithm will search for the adjacent run, R$_2$. After finding it, the algorithm will calculate the power between R$_1$ and R$_2$, which is equal to the height of v$_1$. Given that we only have one power calculated, this will lead to no merges and R$_2$ will be added to our stack. Afterwards, the algorithm will proceed to the right T '' subtree, which is rooted in v$_2$. The algorithm will find R$_3$ and, then, calculate the power between R$_2$ and R$_3$. This will result in the height of v, which is smaller than that of v$_1$, our previously calculated power, leading to the merge of runs R$_1$ and R$_2$ into run R$_{12}$. 

After this, our Forgetful 2-Stack will contain R$_{12}$ and R$_3$. Then, the algorithm will find run R$_4$. Calculating the power between R$_3$ and R$_4$ leads to height v$_2$, which is bigger than that of v, our previously calculated power, meaning the algorithm will not merge any runs and will add R$_4$  to the Forgetful 2-Stack. This will cause the stack to “forget” about R$_{12}$, meaning it will only contain R$_3$ and R$_4$. Following the logic of what the calculation of the power between to runs means, the power between R$_3$ and R$_4$ is higher than that of R$_{12}$ and R$_3$, therefore, the merge between R$_3$ and R$_4$ has to happen before that of R$_{12}$ and R$_3$, so, the algorithm will merge R$_3$ and R$_4$, leading to the processing of both subtrees of T’ and T’’.  

\textit{[Comentário em português: Uma das coisas que ficou aberta no entendimento do algoritmo modificado foi o que fazer quando alcançamos o final do vetor. Acho que a ideia é dar merge para trás até alcançarmos a altura do último merge que aconteceu antes do fim do vetor. No caso de que nenhum merge aconteceu, podemos dar merge para trás subindo C-1 níveis.]}

\textbf{Hypothesis:}
Suppose the proposition holds true for all subtrees of height h $\leq$ k in a powersort merge tree.

\textbf{Inductive Step:} Let T be a Powersort merge tree of height H. We will prove that a (k+2)-Bounded Powersort processes all of T's subtrees that have height of h = k+1. Let T' be one of those subtrees and v its root. We will analyze the sub-trees T$_\ell$ and T$_r$, each rooted in the left and right children of v, v$_\ell$ being the left child and v$_r$ the right child. The height of the T$_\ell$ is, at most, k, therefore, we can apply our induction hypothesis and process T$_\ell$, thus, v$_\ell$ will contain run R$_\ell$, which is the result of the merge of all runs in T$_\ell$. The processing of this sub-tree will happen once the most-left run of T$_r$ is analyzed by the algorithm.\footnote{Once the power between the most-right run of T$_\ell$ with most-left run of T$_r$ is calculated, it will result in the height of v, which is smaller than the height of v$_\ell$, meaning the algorithm will merge everything left in our Forgetful k+2-Stack, for all powers calculated between runs in T$_\ell$ is greater than v$_\ell$} By the end of the processing of T$_\ell$, our Forgetful (k+2)-Stack will contain R$_\ell$ and the most-left run of T$_r$. Given that, by hypothesis, the processing of T$_r$ requires one pass of a (k+1)-Bounded Powersort\footnote{Remember that T$_\ell$ and T$_r$ have, at max, height of k}\footnote{Requiring a (k+1)-Bounded Powersort means the processing needs k+1 free spaces in the Forgetful C-Stack, not a specific tuning of the algorithm.}, and that our Forgetful (k+2)-Stack has exactly k spaces left, with one of them already being occupied by one run of T$_r$, the algorithm will be able to process T$_r$. Similarly to what happened during the processing of T$_\ell$, the calculation of the power between the right-most run of T$_r$ (which is also the right-most run of T') and the next run, which is the left-most run of another subtree of T with height k+1, will result in the height of v - 1.\footnote{This is due to the fact that both runs are in adjacent subtrees: v is the left child of a node v$_p$ and there is a node v' which is the right child of the same v$_p$, the first common ancestor between v and v' is v$_p$, therefore, the first common ancestor between runs in each subtree (the subtree rooted in v and the subtree rooted in v') is v$_p$}. Given that all the powers calculated between runs in T' are greater than v-1, the algorithm will merge all the remaining runs in our Forgetful (k+2) stack, leading to the processing of T$_r$ but also to the processing of T', given that R$_\ell$ will still be on the stack alongside all runs from T$_r$. 
Therefore, if the (k+2)-Bounded Powersort was able to process a generic T' subtree of T with height k+1, it is fair to assume that it will be able to further process all the remaining subtrees of same height by following the same logic as the one described for T'.
\qed

\textit{[Comentário em português: Seria, aqui, necessário, explicitar novamente como funcionaria o algoritmo no final do vetor? Já está no caso base, mas...]}
\newpage
\textbf{Proposition:} It takes C-Bounded Powersort at most $\bigO(n \cdot \log{n})$ element scans to order an array of size n. 
\begin{proof}
Each pass through the array scans element i, with 0 $\leq$ i $\leq$ n, at least once. Given that the algorithm will do $\frac{\log{n}}{c-1}$ passes\footnote{The height of a Powersort merge tree T is $\log n$ and the C-Bounded Powersort
algorithm starts by processing subtrees of T that have height less than or equal to c-1. This results in a smaller tree T' of height $\log n$ - (c-1). A second pass will, once again, process the subtrees of T' that have height less than or equal to c-1, resulting in a smaller tree T''. Therefore, if k is the amount of trees
created by this process until the last tree is of height less than or equal to c-1, meaning it will only require one last pass of the algorithm to be processed, k$\cdot$(c-1) $\leq$ $\log n$}, i will be read
at least $\frac{\log{n}}{c-1}$ times. Once the entire array is ordered, i will also have been read at least twice for each merge it was present in, i.e, 2$\cdot$d$_i$, with d$_i$ being equal to the height of T. Therefore,
the total number of scans is:
\begin{equation}
Scans \leq \sum^{n}_{i=0}2\cdot \log{n} + n \cdot \frac{\log{n}}{c-1}
\end{equation}
\begin{equation}
Scans \leq n\log{n}\cdot(2 + \frac{1}{c-1})
\end{equation}
\begin{equation}
Scans \leq \bigO(n \cdot \log{n})
\end{equation}
\end{proof} 
\end{document}